\documentclass[a4paper, 11pt]{article}

\begin{document}
TODO example figures
TODO which photos did we test on

This project has been quite unpredictable so far. 
Techniques I had hoped would work well have been underwhelming or failed outright.

My initial approach was a simple disparity map.
I hoped to get some intrinsic variables, not only for stereo rectification but also to compute depth in physical units instead of pixels.
I was able to learn the type of camera used for the pictures to get a sensor size, and had a presumed focal length from the lens they used, assuming it was fully zoomed in.
I tried for some time to find a series of shots in the dataset that contained enough of a grid to do some calibration, but to no avail.
Even without camera calibration to calculated specific depths, though, I could still create a disparity map with OpenCV's \verb|stereoRectifyUncalibrated| and interpret the results in pixels to find an area with low change in depth to call the backshore.
However, trying this with SIFT keypoint matching, getting a fundamental matrix with \verb|findFundamentalMat| running RANSAC, provided lacking results.
I spent (wasted) a lot of time tweaking variable, seeing if the apprach was valid but just needed some nudging.
I changed my keypoint match distance ratio, tried various image downscaling factors, modified the RANSAC threshold, and, of course, played with SGBM's \verb|minDisparity|, \verb|numDisparities|, and \verb|blockSize| parameters.

Interestingly, downscaling the image actually helped somewhat, I imagine it got rid of some repetitive patterns in the sand, grasses, and trees.
Matching parameters like match distance ratio and RANSAC threshold didn't do much, the keypoints were mostly good and the rectification process seemed to work really well even without calibration.
It all came down to adjusting the disparity map parameters.
A higher block size was often beneficial. 
Since we're looking for a large pattern from a big feature, we want more signal, less noise.
However, it often just got too big, causing it to never match disparities leading to vertical artifacts.
Setting minimum disparities and number of disparities had little effect.

After a lot of time and effort, I decided to take a break on disparity maps and try a new approach.
This is where my planned timeline really fell apart, since I didn't have a lot of back-up plans in mind.
I turned to more complex structure from motion techniques, since I had a lot of high quality images to work with.
Initially I tried OpenSFM but it was not working so I started using COLMAP and their python bindings PyCOLMAP.
Using eight photos of the dunes near Fort Ord and COLMAP's incremental mapping algorithm, I was able to reconstruct TODO how many points in space.
From these, though, you could only see a general surface, not nearly enough to discern the backshore.

So, I turned instead to their stereo fusion algorithm.
This is where things start to get a little out of my depth but at this point I will do anything to make this work.

\end{document}

